\documentclass[10pt,a4paper]{amsart}
\usepackage{amsmath,amssymb,amsthm,graphicx}
\usepackage{verbatim}
\newenvironment{reference}{\comment}{\endcomment}
%\newenvironment{reference}{}{}
\newenvironment{slogan}{\comment}{\endcomment}
\newenvironment{history}{\comment}{\endcomment}
% For commutative diagrams we use Xy-pic
\usepackage[all]{xy}

% We use 2cell for 2-commutative diagrams.
\xyoption{2cell}
% Theorem environments.
%
\theoremstyle{plain}
\newtheorem{theorem}[subsection]{Theorem}
\newtheorem{proposition}[subsection]{Proposition}
\newtheorem{lemma}[subsection]{Lemma}

\theoremstyle{definition}
\newtheorem{definition}[subsection]{Definition}
\newtheorem{example}[subsection]{Example}
\newtheorem{exercise}[subsection]{Exercise}
\newtheorem{situation}[subsection]{Situation}

\theoremstyle{remark}
\newtheorem{remark}[subsection]{Remark}
\newtheorem{remarks}[subsection]{Remarks}

\numberwithin{equation}{subsection}

% Macros
%
\def\lim{\mathop{\mathrm{lim}}\nolimits}
\def\colim{\mathop{\mathrm{colim}}\nolimits}
\def\Spec{\mathop{\mathrm{Spec}}}
\def\Hom{\mathop{\mathrm{Hom}}\nolimits}
\def\Ext{\mathop{\mathrm{Ext}}\nolimits}
\def\SheafHom{\mathop{\mathcal{H}\!\mathit{om}}\nolimits}
\def\SheafExt{\mathop{\mathcal{E}\!\mathit{xt}}\nolimits}
\def\Sch{\mathit{Sch}}
\def\Mor{\mathop{\mathrm{Mor}}\nolimits}
\def\Ob{\mathop{\mathrm{Ob}}\nolimits}
\def\Sh{\mathop{\mathit{Sh}}\nolimits}
\def\NL{\mathop{N\!L}\nolimits}
\def\CH{\mathop{\mathrm{CH}}\nolimits}
\def\proetale{{pro\text{-}\acute{e}tale}}
\def\etale{{\acute{e}tale}}
\def\QCoh{\mathit{QCoh}}
\def\Ker{\mathop{\mathrm{Ker}}}
\def\Im{\mathop{\mathrm{Im}}}
\def\Coker{\mathop{\mathrm{Coker}}}
\def\Coim{\mathop{\mathrm{Coim}}}

% Boxtimes
%
\DeclareMathSymbol{\boxtimes}{\mathbin}{AMSa}{"02}

%
% Macros for moduli stacks/spaces
%
\def\QCohstack{\mathcal{QC}\!\mathit{oh}}
\def\Cohstack{\mathcal{C}\!\mathit{oh}}
\def\Spacesstack{\mathcal{S}\!\mathit{paces}}
\def\Quotfunctor{\mathrm{Quot}}
\def\Hilbfunctor{\mathrm{Hilb}}
\def\Curvesstack{\mathcal{C}\!\mathit{urves}}
\def\Polarizedstack{\mathcal{P}\!\mathit{olarized}}
\def\Complexesstack{\mathcal{C}\!\mathit{omplexes}}
% \Pic is the operator that assigns to X its picard group, usage \Pic(X)
% \Picardstack_{X/B} denotes the Picard stack of X over B
% \Picardfunctor_{X/B} denotes the Picard functor of X over B
\def\Pic{\mathop{\mathrm{Pic}}\nolimits}
\def\Picardstack{\mathcal{P}\!\mathit{ic}}
\def\Picardfunctor{\mathrm{Pic}}
\def\Deformationcategory{\mathcal{D}\!\mathit{ef}}
\usepackage{cleveref}
\crefname{originallemma}{original lemma}{original lemmas}
\Crefname{originallemma}{Original lemma}{Original lemmas}
\makeatletter
\newlabel{lemma-Noetherian-topology}{{lemma-Noetherian-topology}{source}{}{sourcepointer-1}{}}
\newlabel{lemma-Noetherian-topology@cref}{{[originallemma][0][]lemma-Noetherian-topology}{[1][0][]source}{}{}{}}
\makeatother
\crefname{originallemma}{original lemma}{original lemmas}
\Crefname{originallemma}{Original lemma}{Original lemmas}
\makeatletter
\newlabel{lemma-kolmogorov}{{lemma-kolmogorov}{source}{}{sourcepointer-2}{}}
\newlabel{lemma-kolmogorov@cref}{{[originallemma][0][]lemma-kolmogorov}{[1][0][]source}{}{}{}}
\makeatother
\crefname{originallemma}{original lemma}{original lemmas}
\Crefname{originallemma}{Original lemma}{Original lemmas}
\makeatletter
\newlabel{lemma-locally-Noetherian-quasi-sober}{{lemma-locally-Noetherian-quasi-sober}{source}{}{sourcepointer-3}{}}
\newlabel{lemma-locally-Noetherian-quasi-sober@cref}{{[originallemma][0][]lemma-locally-Noetherian-quasi-sober}{[1][0][]source}{}{}{}}
\makeatother
\crefname{originalreference}{original reference}{original references}
\Crefname{originalreference}{Original reference}{Original references}
\makeatletter
\newlabel{topology-definition-generic-point}{{topology-definition-generic-point}{source}{}{sourcepointer-4}{}}
\newlabel{topology-definition-generic-point@cref}{{[originalreference][0][]topology-definition-generic-point}{[1][0][]source}{}{}{}}
\makeatother
\title{The space of a locally Noetherian decent stack is sober}
\author{The Stacks Project Authors}
\date{stacks-project:stacks-morphisms.tex:local-snapshot-20261012; GNU FDL 1.2 or later}
\begin{document}
\maketitle
\noindent\textit{Editorial note.} All delivered mathematical text is the author's own, copied from the pinned Stacks Project chapter source sources/stacks-algebraic-stacks-morphisms/pinned-original-source.tex; no statement or proof step was rewritten, reordered or removed.\par\medskip
\noindent\textit{Editorial note.} Preamble provenance: the Stacks Project keeps one single shared preamble for every chapter; this card takes it directly from the staged sources/stacks-algebraic-stacks-morphisms/preamble.tex as three verbatim spans (lines 8-12, 14-18 and 154-217), so no derived file was written and nothing was added or redefined. The wrapper class and the base packages are supplied by the assembler.\par\medskip
\noindent\textit{Editorial note.} Bibliography provenance: neither the counted statement nor its proof contains a citation (this lane rescans the whole span before writing the card), so the card carries no bibliography and no entry was reconstructed.\par\medskip
\noindent\textit{Editorial note.} Selection provenance: this card was opened under the file-based control contract (CONTROL.md), with the mechanical pre-card gate run on its statement and proof interval before the card was written.\par\medskip
\noindent\textit{Allowed standard prerequisites (editorial).}\begin{itemize}
\item The statement that the space of a locally Noetherian algebraic stack is locally Noetherian. Invoked at the start of the counted proof (source lines 10565-10567). Supplies the first ingredient.
\item The statement that the space of a decent stack is Kolmogorov. Invoked inside the counted proof (source lines 10568-10569). Supplies the second ingredient.
\item The statement that a locally Noetherian Kolmogorov space is quasi-sober. Invoked near the end of the counted proof (source lines 10570-10571). Supplies the third ingredient.
\item The definition of sobriety through existence of generic points. Invoked at the end of the counted proof (source line 10572) as a cross-chapter pointer. Assembles the three ingredients into the conclusion.
\end{itemize}\medskip
\ifdefined\hypertarget\hypertarget{sourcepointer-1}{}\fi
\noindent\textit{Source reference.} Original lemma lemma-Noetherian-topology: sources/stacks-algebraic-stacks-morphisms/pinned-original-source.tex, source lines 2253--2253. Navigation pointer only: the counted text refers to the Stacks item lemma-Noetherian-topology, which is an existing labeled result of this same chapter and is not part of this card. The pointer carries the target's own label as its number and its definition line as its span; its content is not reproduced here.\par\medskip
\ifdefined\hypertarget\hypertarget{sourcepointer-2}{}\fi
\noindent\textit{Source reference.} Original lemma lemma-kolmogorov: sources/stacks-algebraic-stacks-morphisms/pinned-original-source.tex, source lines 10519--10519. Navigation pointer only: the counted text refers to the Stacks item lemma-kolmogorov, which is an existing labeled result of this same chapter and is not part of this card. The pointer carries the target's own label as its number and its definition line as its span; its content is not reproduced here.\par\medskip
\ifdefined\hypertarget\hypertarget{sourcepointer-3}{}\fi
\noindent\textit{Source reference.} Original lemma lemma-locally-Noetherian-quasi-sober: sources/stacks-algebraic-stacks-morphisms/pinned-original-source.tex, source lines 2280--2280. Navigation pointer only: the counted text refers to the Stacks item lemma-locally-Noetherian-quasi-sober, which is an existing labeled result of this same chapter and is not part of this card. The pointer carries the target's own label as its number and its definition line as its span; its content is not reproduced here.\par\medskip
\ifdefined\hypertarget\hypertarget{sourcepointer-4}{}\fi
\noindent\textit{Source reference.} Original reference topology-definition-generic-point: sources/stacks-algebraic-stacks-morphisms/pinned-original-source.tex, source lines 2283--2283. Navigation pointer only: the counted text refers to the Stacks item topology-definition-generic-point, which belongs to another chapter of the project and is not part of this card. The pointer carries the target's own label as its number because this chapter's staged snapshot does not contain the other chapter's tag table; its content is not reproduced here.\par\medskip
\begin{lemma}
\label{lemma-decent-locally-noetherian-sober}
Let $\mathcal{X}$ be a decent, locally Noetherian algebraic stack.
Then $|\mathcal{X}|$ is a sober locally Noetherian topological space.
\end{lemma}

\begin{proof}
By Lemma \ref{lemma-Noetherian-topology} the topological space $|\mathcal{X}|$
is locally Noetherian. By Lemma \ref{lemma-kolmogorov} the topological space
$|\mathcal{X}|$ is Kolmogorov. By
Lemma \ref{lemma-locally-Noetherian-quasi-sober} the topological space
$|\mathcal{X}|$ is quasi-sober. This finishes the proof, see
Topology, Definition \ref{topology-definition-generic-point}.
\end{proof}

\end{document}
